% ============================================================
%  Introduction Générale
%  Corps : 12pt Times New Roman, interligne 1.5
% ============================================================
\clearpage


\vspace*{2cm}

\begin{center}
{\fontsize{24}{29}\selectfont\bfseries Introduction Générale}
\end{center}
\addcontentsline{toc}{chapter}{Introduction Générale}

\vspace{1.5cm}

La gestion d'un cabinet dentaire implique un ensemble de processus complexes~: la prise de rendez-vous, la gestion des dossiers patients, le suivi des praticiens et l'administration des créneaux horaires. Les systèmes traditionnels, souvent basés sur des registres papier ou des tableurs, présentent des limitations significatives en termes de fiabilité, de traçabilité et d'efficacité.

Dans un cabinet dentaire, la planification des rendez-vous est une tâche critique. Une mauvaise gestion peut entraîner des conflits de créneaux (deux patients programmés chez le même médecin à la même heure), des temps d'attente excessifs, des oublis de rendez-vous, et une sous-utilisation des ressources (cabinets vides). Ces problématiques affectent directement la qualité du service et la satisfaction des patients.

Par ailleurs, les dossiers patients --- contenant les informations personnelles, les antécédents médicaux, les contrats d'assurance et l'historique des consultations --- doivent être centralisés, sécurisés et facilement accessibles pour permettre un suivi efficace.

Dans ce contexte, le présent projet vise à concevoir et réaliser une application web moderne baptisée \textbf{Dental Reservation}, capable de répondre à l'ensemble de ces besoins. L'application offre une interface intuitive permettant aux praticiens de gérer leurs patients, planifier et modifier les rendez-vous avec détection automatique des conflits, visualiser les créneaux disponibles sur un calendrier hebdomadaire, et consulter les dossiers médicaux --- le tout de manière sécurisée grâce à un mécanisme d'authentification par jeton JWT.

Ce rapport décrit le projet dans son intégralité~:

\begin{itemize}
    \item \textbf{Chapitre 1} présente le contexte général, l'organisme d'accueil et un aperçu de la solution.
    \item \textbf{Chapitre 2} détaille l'étude fonctionnelle et technique~: architecture, besoins et spécifications.
    \item \textbf{Chapitre 3} couvre la conception UML avec les diagrammes de cas d'utilisation, de classes, de séquence et d'activité.
    \item \textbf{Chapitre 4} présente la réalisation, avec une description détaillée de chaque module et fonctionnalité de l'application.
\end{itemize}

\clearpage
